\documentclass[12pt]{article}

\usepackage[margin=1in]{geometry}
\usepackage{setspace}
\usepackage{booktabs}
\usepackage{graphicx}
\usepackage{hyperref}
\usepackage{enumitem}
\usepackage{amsmath}

\onehalfspacing

\title{Climate Risk, Resilience Investment, and Municipal Bond Pricing in New Jersey}
\author{Rutgers Economic Advisory Laboratory}
\date{\today}

\begin{document}
\maketitle

\begin{abstract}
Briefly summarize the research question, data sources, empirical approach, headline results, and policy relevance.
This paragraph should be written last.
It should clearly state that the current evidence is based on an exploratory secondary-market bond panel and a synthetic AAA benchmark unless final validated benchmark and issuance data are added before submission.
\end{abstract}

\section{Introduction}
Explain why climate physical risk and adaptation investment may matter for municipal borrowing costs.
Introduce the New Jersey policy context, including coastal flood exposure, municipal infrastructure finance, and NJDEP/NJ I-Bank resilience programs.
State the core research question: whether exposed municipalities face higher bond spreads and whether resilience participation is associated with lower or higher spreads.
End with a concise preview of the paper's contribution and corrected empirical takeaway.

\section{Policy and Institutional Background}
Describe the relevant New Jersey climate adaptation and financing institutions.
Include NJDEP, NJ I-Bank, Water Bank/PPL processes, CHAMP, FEMA CRS, MS4 stormwater context, and the 2022 policy pivot if it remains central to the analysis.
Clarify what the 150-point resilience criterion means, which fiscal years it applies to, and how it should affect municipal incentives in theory.

\section{Conceptual Framework}
Lay out the mechanisms through which climate risk and resilience may affect borrowing costs.
Discuss expected signs for:
\begin{itemize}[leftmargin=*]
    \item physical-risk penalties from sea-level-rise exposure;
    \item resilience premiums if adaptation reduces expected losses;
    \item disclosure or signaling penalties if program participation reveals risk;
    \item capital-expenditure or liquidity concerns if adaptation requires near-term borrowing.
\end{itemize}
This section should avoid claiming any mechanism is proven unless supported by the final empirical results.

\section{Data}
Describe each dataset, its source, time coverage, unit of observation, and cleaning steps.
Subsections should include the following.

\subsection{Municipal Boundary and Climate Exposure Data}
Explain how municipal boundaries are joined to flood or sea-level-rise exposure measures.
Define \texttt{slr\_exposure\_pct} as the percent of total municipal market value exposed under the 4-foot SLR scenario from \texttt{data/data\_cleaned/climate/nj\_statewide\_frm\_scores\_FINAL.csv}.
Discuss why tax-base-weighted exposure is preferable to raw land-area exposure, while noting uncertainty in the flood and parcel valuation layers.

\subsection{Bond Market Data}
Describe the WRDS secondary-market municipal trade data, the GO-like trade filter, municipality matching, and the 2015--2025 coverage.
State clearly whether the dependent variable uses secondary trades or issuance/offering yields.
If the synthetic AAA benchmark remains in use, explain how it is constructed and label it as a proxy rather than Bloomberg/MMD AAA.
If true MMD data are later added, describe the same-date, same-maturity spread construction here.

\subsection{Resilience and Adaptation Data}
Explain the available resilience indicators: CHAMP applicant PDFs, FEMA CRS participation/class, MS4 fields, and any Water Bank/PPL information used.
Define \texttt{ever\_champ}, \texttt{first\_champ\_sfy}, \texttt{best\_champ\_rank}, \texttt{crs\_resilient}, and the time-varying \texttt{is\_resilient} flag.
Be explicit that the CHAMP flag is based on locally parsed applicant PDFs unless a fully validated I-Bank treatment history is added.

\subsection{Fiscal and Demographic Controls}
Describe UFB debt-to-assessed-value controls and ACS median household income controls.
Note any forward filling, missing data rules, or years with incomplete coverage.

\section{Sample Construction and Entity Resolution}
Document the full pipeline from raw files to \texttt{final\_panel\_master.csv}.
Report the number of trade rows, municipalities, years, treated/cohort municipalities, and dropped observations.
Describe how municipality names are normalized and how ambiguous municipality names are handled.
Include a table showing sample counts by year and treatment/cohort status.

\section{Empirical Strategy}
Present the main regression models in equations.
Define the pooled OLS, two-way fixed effects DiD, triple-difference specification, and event-study/pre-trend check.
For example:
\[
Spread_{i t} = \alpha_i + \lambda_t + \beta_1(Champ_i \times Post_t)
 + \beta_2(Post_t \times SLR_i)
 + \beta_3(Champ_i \times Post_t \times SLR_i)
 + X_{it}\gamma + \epsilon_{it}.
\]
Explain what each coefficient measures and what assumptions are required for causal interpretation.
State that standard errors are clustered by municipality.
State clearly which controls enter each model: the current preferred Model B and Model C specifications include time to maturity but intentionally exclude debt-to-AV.
Discuss why the current specification is exploratory if secondary trades, synthetic benchmarks, and missing bond-level controls remain.

\section{Data Quality and Validation}
Summarize validation checks that should be shown before final submission.
Include checks for missing municipalities, unmatched trade descriptions, outlier spreads, duplicate CUSIP-date observations, treatment timing, and pre/post support by municipality.
Discuss winsorization of spreads at the 1st and 99th percentiles and justify whether it is retained in final specifications.

\section{Results}
Report the main coefficient table and interpret magnitudes in basis points.
Include subsections for:
\begin{itemize}[leftmargin=*]
    \item baseline climate exposure and CHAMP cohort associations;
    \item two-way fixed effects DiD results;
    \item coastal triple-difference results;
    \item pre-trend or event-study evidence.
\end{itemize}
Based on the current corrected results, state that the triple interaction is not statistically distinguishable from zero and that the pre-trend check weakens causal interpretation.
Update this section if additional validated data change the estimates.

\section{Robustness Checks}
List robustness checks to include in the final paper.
Recommended checks include raw versus winsorized spreads, excluding Atlantic City or other distressed outliers, restricting to towns with balanced pre/post observations, alternative SLR thresholds, CRS-based resilience definitions, CHAMP-only definitions, and alternative clustering or weighting choices.
If true MMD benchmarks are added, compare synthetic-proxy results with MMD-based results.

\section{Discussion}
Interpret what the evidence does and does not imply for municipal bond markets.
Discuss possible explanations for null, positive, or negative estimates without overclaiming.
Relate findings back to the conceptual framework and policy context.
Clearly distinguish market-pricing evidence from broader social benefits of resilience investments.

\section{Policy Implications}
Translate the findings for NJDEP, NJ I-Bank, municipal finance officers, and local planners.
If results remain null or exploratory, emphasize data infrastructure, disclosure quality, treatment tracking, and the need for better bond-level controls rather than recommending subsidies based on an unproven penalty.
If final validated results show a robust penalty or premium, explain the policy implications and fiscal magnitude carefully.

\section{Limitations}
Be direct about remaining constraints.
Current limitations include the synthetic AAA benchmark, secondary-market rather than issuance data, missing rating/insurance/tax/call/issue-size controls, incomplete treatment validation, missing 2020 bond data, and partial municipality coverage.
Explain how each limitation could bias results or limit interpretation.

\section{Conclusion}
Restate the research question, summarize the most defensible empirical finding, and identify the next data improvements needed for a final policy-grade estimate.
Avoid introducing new evidence in this section.

\section*{References}
Add citations for NJDEP climate reports, NJ PACT/SLR documentation, FEMA CRS documentation, MSRB/EMMA or WRDS data documentation, MMD/Bloomberg benchmark documentation if used, and relevant academic municipal bond climate-risk literature.

\appendix

\section{Variable Definitions}
Create a table defining every variable used in the analysis, including source, unit, expected sign, and construction notes.

\section{Data Pipeline Reproducibility}
List the commands needed to reproduce the cleaned panel and results:
\begin{verbatim}
make install
make process-all
python3 data_processing/run_did_regression.py
python3 data_processing/export_did_coefficients.py
\end{verbatim}
Note any dependencies or proprietary data files required outside the repository.

\section{Additional Tables and Figures}
Reserve space for descriptive statistics, sample coverage maps, coefficient tables, event-study plots, and robustness tables.

\end{document}
